\section{What v1.0 would have to show}
\label{sec:towardone}

Two sentences a future release might want to write are ``FishBot v1.0 is near-optimal'' and ``FishBot
v1.0 is the strongest known Fish-playing team''. Neither is supportable by anything in this corpus
today, and neither is close. This section states what evidence each would require, so that the gap is
a list of experiments rather than a matter of taste. Nothing here has been measured; every item is
open.

\subsection{For ``near-optimal''}
\label{sec:towardone-optimal}

Near-optimality is a claim about distance to a benchmark, so the first requirement is a benchmark.
The corpus has none.

\begin{enumerate}

\item \textbf{An upper bound on exploitability, not a lower one.} Every exploitability figure in this
  corpus, including all eight arms of \S\ref{sec:exploit}, is a lower bound produced by a search: it
  says what \emph{this} search found, and computing the true value is NP-hard even white-box. A
  near-optimality claim needs a bound in the other direction --- an argument or a computation showing
  that no adversary in a stated class can exceed some $\varepsilon$. No such bound exists here for any
  class, and producing one for even a restricted class would be the single most valuable result the
  next cycle could return.

\item \textbf{A team analogue of local best response, and the rollout policy it needs.} The
  poker line's standard exploitability probe rests on a cheap, correct-enough passive continuation ---
  ``what if nobody does anything clever from here''. Fish has no passive continuation: a player must
  ask or declare. Until a defensible rollout policy for that role is named and calibrated, the probe
  cannot be ported, and the literature's own warning that choosing a value function is the hard part
  applies directly.

\item \textbf{An adversary that is genuinely ex-ante correlated.} The threat model names three
  correlation regimes and makes the correlated one the headline. What has been built and measured is
  the synchronized regime, plus one correlated fit that reached \vsevenCorrelatedFit{} against a class
  floor of \vsevenCorrelatedFloor{} --- i.e. nothing. Whether a genuinely correlated three-seat
  responder is buildable at all for a game this size is open, and the honest position is that the
  headline regime has fallen back to a weaker one.

\item \textbf{A detection floor low enough for the claim.} The instrument's floor is
  \vsevenPrimaryFloor{} points and \S\ref{sec:corrections-corpus} records that it does not buy down
  with evaluation games. ``Near-optimal'' is not a statement compatible with an instrument that cannot
  distinguish a policy from one a point and a half weaker. Lowering the floor is not an evaluation-power
  problem --- it is a search-power problem, and phase~2 established that by measurement.

\item \textbf{The sum of small effects, addressed.} This report and its predecessor both conclude that
  no \emph{individual} mechanism on the register is large. Neither addresses whether a dozen
  quarter-point mechanisms exist and compound. Resolving a quarter of a point on the scoreboard needs
  roughly \vsevenScoreboardGames{} games a cell; the per-decision estimator of
  \S\ref{sec:rejected-perdecision} does it in
  \vsevenPerDecisionGamesLo{}--\vsevenPerDecisionGamesHi{}. That estimator is confirmed and unused,
  because fitting against it fails. Using it as an \emph{estimator} while fitting against the
  scoreboard is the obvious unattempted experiment.

\end{enumerate}

\subsection{For ``the strongest known Fish-playing team''}
\label{sec:towardone-strongest}

This is a comparative claim and it is currently unfalsifiable, because the comparison class has one
member.

\begin{enumerate}

\item \textbf{An opponent this project did not write.} Every arm, archetype and adversary in this
  corpus was built inside it. The panel's deception archetypes are stylised implementations of
  manoeuvres one person reported from live play. A claim to be strongest known requires at least one
  independently authored opponent, and ideally a published rule dialect other people can implement
  against.

\item \textbf{Play against people, measured.} The corpus has no human-play measurement of any kind.
  What it has is a partner-regime table in which ``a foreign partner'' is another bot
  (\S\ref{sec:results-partners}). The project owner's standing decision --- never headline the
  self-play configuration as the one a person will meet --- exists precisely because the substitute is
  not the thing. A strongest-known claim needs a duplicate-style human battery with enough deals to
  resolve a point, and the six-seat duplicate design for it has not been settled.

\item \textbf{A stated, published dialect.} Fish is played in many house dialects, and
  \S\ref{sec:results-dialects} shows the arm ordering is stable across eight of them --- which is
  evidence the claim could survive dialect variation, and also evidence that ``Fish'' names a family
  rather than a game. Any comparative claim has to name the dialect it holds in.

\item \textbf{Somebody else's evaluation.} Every number in this report was produced by one engine,
  written by this project, audited by tests this project wrote. The strongest available substitute for
  external replication was used --- material sealed before the configuration existed, digests computed
  by a second build, a protocol committed in advance --- and it is not the same thing.

\end{enumerate}

\subsection{The next cycle's first three jobs}
\label{sec:towardone-next}

Stated concretely, because a barrier report should hand over work rather than a mood.

\begin{enumerate}
\item \textbf{Fix the cross-deal agent residue, and re-measure everything the fix moves.} The carrier
  is named (\texttt{\vsevenSsixCarrier{}}); the member is not. The repair changes the play of every
  searching configuration by about one deal in \vsevenSsixIncidence{}. It is the first maintenance
  item after this freeze, and it must precede the next freeze rather than follow it.
\item \textbf{Test whether the small effects compound.} Per-decision estimation, per-game objective,
  on a lattice of mechanisms each individually below the floor. This is the one hypothesis both this
  cycle and the last explicitly declined to test, and both said so.
\item \textbf{Build an upper bound for one restricted adversary class.} Even a loose bound over a
  small class would convert every exploitability sentence in this corpus from ``we did not find one''
  to ``there is not one larger than $\varepsilon$'', which is a different kind of claim and the only
  kind that supports the word near-optimal.
\end{enumerate}
